\section{Recursive words with incompatible displacements}
\label{sec:construction}

We now build the words to which Proposition~\ref{prop:displacement}
will be applied. The row list must be long enough that its one-row shift
remains inexpensive after multiplication by $(2P)^{2h}$ in that
inequality. At the same time, it must be shorter than the
product of many component periods, so that a row shift cannot cancel
too many prime-period differences. The following choices provide both
properties.

Fix a sufficiently large integer $k$ and set
\begin{equation}\label{eq:parameters}
 b=100k,\qquad P=b^2,\qquad h=b/10=10k,\qquad L=P^{b/2}.
\end{equation}
Choose $b$ distinct primes in $[P,2P]$, and call their set
$\mathcal P$. The prime number theorem gives
$\pi(2P)-\pi(P)\sim P/\log P>b$ for all sufficiently large $k$;
the explicit interval estimate of Rosser and Schoenfeld
\cite[Corollary~3, equation~(3.8)]{RosserSchoenfeld1962} also gives
this consequence. We use this one pool and the same parameters at
every level.

\subsection{The recursive family and its cheap steps}

For every $0\le s\le k$ and $p\in\mathcal P$, the construction supplies
a finite abelian state group $G_{s,p}$, a distinguished element
$v_{s,p}$ of order $p$, and a word map
\[
 W_{s,p}:G_{s,p}\longrightarrow\mathcal A_{s,p}^{\,n_s}.
\]
It also supplies a finite list $\mathcal S_{s,p}$ of pairs $(e,r)$,
where $e\in G_{s,p}$ is a step and $r>0$ is its displacement budget.
The word map itself need not be injective. Define
\begin{equation}\label{eq:recurrences}
 \alpha_s=24^{-s},\qquad Q_0=2,\qquad
 Q_s=\frac{2(2P)^{2h}}{L}+\frac{Q_{s-1}}h.
\end{equation}

\begin{proposition}[Recursive properties]\label{prop:family}
The families can be chosen with
\[
 n_s=\bigl(2(2b-1)L\bigr)^s,\qquad
 |\mathcal A_{s,p}|\le A_s,\qquad
 A_0=2P,\quad A_s=1+(2b-1)A_{s-1},
\]
so that the following assertions hold.
\begin{enumerate}
\item Every listed step is cheap at every state:
\begin{equation}\label{eq:step-property}
 \ID\bigl(W_{s,p}(z),W_{s,p}(z+e)\bigr)\le n_s r
 \quad\bigl((e,r)\in\mathcal S_{s,p},\ z\in G_{s,p}\bigr).
\end{equation}
\item Every nonzero distinguished shift is separated at every state:
\begin{equation}\label{eq:far-property}
 \ID\bigl(W_{s,p}(z),W_{s,p}(z+uv_{s,p})\bigr)\ge n_s\alpha_s
 \quad(z\in G_{s,p},\ 1\le u<p).
\end{equation}
\item If $H:G_{s,p}\to\ell_1$ satisfies the averaged budgets
$V_H(e)\le r$ for all $(e,r)\in\mathcal S_{s,p}$, then
\begin{equation}\label{eq:l1-property}
 V_H(uv_{s,p})\le Q_s\qquad(1\le u<p).
\end{equation}
\end{enumerate}
\end{proposition}

The first two assertions describe the geometry of the words. The third
says that any $\ell_1$ map respecting the same cheap-step budgets has
small average displacement in the separated direction. We prove the
proposition by simultaneous induction on $s$, constructing the word maps
and the cheap-step lists together.

At level zero, take
\[
 G_{0,p}=\Z/p\Z,\quad v_{0,p}=1,\quad
 W_{0,p}(z)=(z),\quad
 \mathcal S_{0,p}=\{(e,2):e\ne0\}.
\]
The single symbols are distinct. Their insertion--deletion distance is
two, so the first two assertions hold with $n_0=1$ and $\alpha_0=1$.
The listed budgets give the third assertion with $Q_0=2$.
The alphabet has $p\le2P$ symbols.

Suppose the families at level $s-1$ have been constructed. To build
the family indexed by $p$, enumerate the pool as $p_1,\ldots,p_b$ with
$p_1=p$, and abbreviate
\[
 G_j=G_{s-1,p_j},\quad v_j=v_{s-1,p_j},\quad
 W_j=W_{s-1,p_j},\quad m=n_{s-1}.
\]
Set
$G_{s,p}=\prod_{j=1}^bG_j$ and
$v_{s,p}=(v_1,0,\ldots,0)$, which has order $p$.

For $z=(z_1,\ldots,z_b)$, the new word is a concatenation of $L$ rows,
indexed by $i=0,\ldots,L-1$. Row $i$ begins with a payload of $2b-1$
slots. Its first $b$ slots contain separately tagged copies of
$W_1(z_1+iv_1)$; its remaining slots contain tagged copies of
$W_j(z_j+iv_j)$ for $j=2,\ldots,b$, in that order. Each slot has its
own tag, fixed across all rows and different from every other slot tag.
The payload length is
$B=(2b-1)m$. Append a marker $\#^B$, where $\#$ is absent from all
payloads.

More formally, the new alphabet is the disjoint union of $\{\#\}$ and
one set $\{r\}\times\mathcal A_{s-1,p_j}$ for each slot $r$ carrying
component $j$. This also tags every lower-level marker, so none becomes
the new symbol $\#$. The full length and alphabet bounds are
\begin{equation}\label{eq:row-length}
 n_s=2BL=2(2b-1)L\,n_{s-1},
 \qquad |\mathcal A_{s,p}|\le1+(2b-1)A_{s-1}.
\end{equation}

There are two kinds of cheap steps. First include the simultaneous
step $\tau=(v_1,\ldots,v_b)$ with budget $2/L$. Row $i$ of the word at
$z+\tau$ is row $i+1$ of the word at $z$ whenever $0\le i<L-1$.
Deleting the first row and inserting one final row costs at most
$4B=2n_s/L$.

For the other steps, let $c_1=b$ and $c_j=1$ for $j\ge2$, and put
\begin{equation}\label{eq:weights}
 \omega_j=\frac{c_jm}{2B},\qquad \sum_{j=1}^b\omega_j=\frac12.
\end{equation}
The number $\omega_j$ is the fraction of the full word occupied by
component $j$; markers occupy the other half. For each
$(e,r)\in\mathcal S_{s-1,p_j}$ include its lift $\widetilde e$ to
coordinate $j$ with budget $\omega_jr$. The preceding level's bound
holds at every translated state $z_j+iv_j$. Editing the $Lc_j$ affected
slots separately therefore costs at most
$Lc_jmr=n_s\omega_jr$. This proves \eqref{eq:step-property}.

\subsection{Separation across row boundaries}\label{sec:separation}

We first isolate the matching argument supplied by the fresh markers.
It will convert a comparison of every pair of payloads into a comparison
of the full row lists.

\begin{lemma}[Marked payloads]\label{lem:marked}
Let $T,p,q\ge1$ be integers, let $\Gamma$ be an alphabet, and let
$\#\notin\Gamma$. For $U_a,V_b\in\Gamma^p$ with $1\le a,b\le T$, put
\[
 X=U_1\#^q\cdots U_T\#^q,\qquad
 Y=V_1\#^q\cdots V_T\#^q,\qquad N=T(p+q).
\]
If $d\ge0$ and $r_p(U_a,V_b)\ge d$ for every $a,b$, then
\begin{equation}\label{eq:marked}
 r_N(X,Y)\ge\frac{Tdq}{d+q}\ge\frac T2\min(d,q).
\end{equation}
In particular, $d\le q$ implies $\ID(X,Y)\ge Td$.
\end{lemma}

\begin{proof}
Fix an increasing matching, with $E$ unmatched positions on each side.
No source payload symbol can match a target marker. Let $s$ be the
number of source payloads whose matches reach at least two target
payloads.

A source payload whose matches lie in one target payload loses at least
$d$ source positions. A payload with no matches loses $p\ge d$ positions.
Summing over these $T-s$ payloads gives $E\ge d(T-s)$.

For each of the other $s$ payloads, choose a target marker between two
of its matches, as in Figure~\ref{fig:marker}. The entire marker is
unmatched: a matched target
position between the two chosen matches would have its source partner
between their source partners, inside this source payload, where $\#$
does not occur. The target matched spans of different source payloads
are ordered and disjoint, so the chosen markers are distinct. Thus
$E\ge sq$.

Combining the two estimates gives $E\ge d(T-E/q)$ and hence
$E\ge Tdq/(d+q)$. The latter is at least
$(T/2)\min(d,q)$. This holds for every matching, including when $d=0$,
and proves \eqref{eq:marked}. Finally $\ID(X,Y)=2r_N(X,Y)$.
\end{proof}

\begin{figure}[t]
\centering
\begin{tikzpicture}[x=0.9cm,y=0.9cm,>=Latex,font=\small]
 \draw[fill=blue!6] (1.5,1.9) rectangle (8.5,2.5);
 \node at (5,2.2) {source payload (no $\#$)};
 \draw[fill=blue!6] (0,0) rectangle (3,0.6);
 \draw[fill=gray!20] (3,0) rectangle (5.3,0.6);
 \draw[fill=blue!6] (5.3,0) rectangle (10,0.6);
 \node at (1.5,0.3) {target payload};
 \node at (4.15,0.3) {$\#^q$};
 \node at (7.65,0.3) {target payload};
 \draw[->,thick] (2.2,1.9)--(2.2,0.67);
 \draw[->,thick] (7.5,1.9)--(7.5,0.67);
 \node[align=center] at (4.8,1.25)
       {no source $\#$\\between these matches};
\end{tikzpicture}
\caption{Two matches from one source payload enclose an unmatched target
marker. The arrows represent matched pairs; no alignment of other payload
boundaries is assumed.}\label{fig:marker}
\end{figure}

We apply the lemma to the words at $z$ and $z+uv_{s,p}$, where
$1\le u<p_1$. Compare source payload $i$ with target payload $i'$ and
write $\delta=i'-i$. Their relative shifts are $(\delta+u)v_1$ in
component $1$ and $\delta v_j$ in component $j\ge2$.

At least $b/3$ slots have a nonzero shift. If $p_1\nmid\delta+u$, all $b$
copies of component $1$ do. Otherwise $\delta\ne0$, since $1\le u<p_1$.
Now $0<|\delta|<L=P^{b/2}$, whereas the product of $b/2$ distinct primes
from the pool is at least $P^{b/2}$. Fewer than $b/2$ pool primes can
therefore divide $\delta$. Among components $2,\ldots,b$, at least
$b-1-b/2\ge b/3$ have a nonzero shift.

In each such slot the induction hypothesis gives
insertion--deletion distance at least $m\alpha_{s-1}$. A matching of
the two complete length-$m$ slot words therefore leaves at least
$m\alpha_{s-1}/2$ source positions unmatched. Since the slot tags are
disjoint, every matching between the two payloads restricts to a
matching of the corresponding complete slot words. Its size is bounded
by their $\LCS$ even if other source payloads use some positions of the
same target slot. Thus every pair of payloads has one-sided deficit at
least
\[
 d_s=\frac{\alpha_{s-1}bm}{6}.
\]
We have $d_s\le B$. Lemma~\ref{lem:marked}, with $T=L$ and marker
length $B$, gives
$\ID(W_{s,p}(z),W_{s,p}(z+uv_{s,p}))\ge Ld_s$. Dividing by
\eqref{eq:row-length},
\[
 \frac{\ID(W_{s,p}(z),W_{s,p}(z+uv_{s,p}))}{n_s}
 \ge\frac{\alpha_{s-1}b}{12(2b-1)}
 \ge\frac{\alpha_{s-1}}{24}=\alpha_s.
\]
This proves the pointwise separation \eqref{eq:far-property}.

\subsection{The averaged analytic recurrence}\label{sec:averaged}

Suppose $H:G_{s,p}\to\ell_1$ meets the listed averaged budgets.
We need the preceding level's assertion for one component, but its
hypotheses concern averages over that component's entire state group.
The right map is a direct sum over all settings of the other components.

For fixed $j$, let $Y_j=\prod_{\ell\ne j}G_\ell$ and define
\begin{equation}\label{eq:slice-map}
 T_j(x)=\frac1{\omega_j|Y_j|}
          \bigoplus_{y\in Y_j}H(x,y),\qquad x\in G_j.
\end{equation}
Here $H(x,y)$ means that $x$ occupies coordinate $j$, and the finite
direct sum of sequence spaces is identified with $\ell_1$. For every
$e\in G_j$, additivity of the sum norm gives
\begin{equation}\label{eq:slice-identity}
 V_{T_j}(e)
 =\frac1{\omega_j}\E_{x\in G_j,\,y\in Y_j}
       \norm{H(x+e,y)-H(x,y)}
 =\frac{V_H(\widetilde e)}{\omega_j}.
\end{equation}
The lifted step budgets therefore make $T_j$ satisfy all the
preceding level's hypotheses at once. The induction assertion yields
\begin{equation}\label{eq:component-control}
 V_H(av_j)\le\omega_jQ_{s-1}\qquad(0\le a<p_j),
\end{equation}
where $a=0$ is immediate. No individual slice of $H$ is required to
satisfy those hypotheses.

Apply Proposition~\ref{prop:displacement} to the factors $G_j$ and
their distinguished elements. The simultaneous budget, component
control and \eqref{eq:weights} give
\begin{align*}
 V_H(uv_{s,p})
 &\le(2P)^{2h}V_H(\tau)
      +\frac2h\sum_{j=1}^b\E_{a\in\Z/p_j\Z}V_H(av_j)\\
 &\le\frac{2(2P)^{2h}}L
      +\frac2h\sum_{j=1}^b\omega_jQ_{s-1}\\
 &=\frac{2(2P)^{2h}}L+\frac{Q_{s-1}}h=Q_s.
\end{align*}
This proves \eqref{eq:l1-property} and completes the simultaneous
induction in Proposition~\ref{prop:family}.

The quotient $\alpha_s/Q_s$ measures the resulting obstruction to
embedding the word image. Indeed, let
$f:W_{s,p}(G_{s,p})\to\ell_1$ have distortion $D$, normalized so that
$\ED(x,y)/D\le\norm{f(x)-f(y)}\le\ED(x,y)$.
Then $H=f\circ W_{s,p}/n_s$ satisfies the listed budgets by
$\ED\le\ID$ and \eqref{eq:step-property}. On the other hand,
\eqref{eq:far-property} and $\ED\ge\ID/2$ give
\[
 \frac{\alpha_s}{2D}\le V_H(v_{s,p})\le Q_s,
 \qquad D\ge\frac{\alpha_s}{2Q_s}.
\]
Thus pointwise word separation can be compared directly with averaged
$\ell_1$ displacement. It remains to estimate this quotient and retain
it when the growing alphabet is replaced by binary symbols.
